\section{The uniform rank bound on every ramified row}

% BEGIN U81L1
\paragraph{U81L1. At most six central generators per unit of central multiplicity.}
Take the physical and coefficient quotients modulo $vN^\Delta$ and
$v\mathcal N$ respectively, restricted to invariant classes at
depth one. V81P2 shows that $\alpha,\beta$ of P81F5 give
well-defined linear maps between these quotients. In each direction
one lifts any representative, using the corresponding inclusion
for $v$ times the integral lattice in V81P2. Their composites
multiply by $p^{2\sigma}$ since the errors $J_P,J_L$ lie inside
the respective $v$-lattices. The maps commute with the row map $F$.
Work at a single chosen letter.

For any $w$ at depth $M$ on the target pass to its depth-one
class. The first frame proposition V81P3 gives an equality in
the coefficient quotient expressing $\beta(\Pi w)$ as a sum of
coordinate tensors and a remainder. Here $\Pi=\pi_{\Omega,1}$,
the coordinate tensors satisfy the individual lifting hypotheses,
and all remainders lie in the image of the \emph{same}
$\mathrm{IIrest}(S_r)$ with $\operatorname{rank} S_r\le2m$
from V81P1. Choose an integral basis of that quotient once for
the whole target letter. Its size is at most $6m$ by P81F2.
Apply $\alpha$. For each basis class write $y_b$ for its physical
image at depth one. Multiplying the equality by
$p^\kappa/\Pi\in O$ (P81F3) gives $p^{2\sigma+\kappa}w$ as a
sum of physical coordinate classes and $O$-multiples of the
$y_b$. Scalars in this equality may be left outside each
coordinate lift.

At $q$ the two branches have the same deep sign and V81P4 lifts
each coordinate after multiplication by $p^\rho$. At $d$ the
two signs are opposite and X81M2--X81M3 give the exponent
$\rho+\max(2\sigma,y_0)$ for every relative position of the
two endpoints. This also dominates the pure-sign exponent.
Consequently for the native
$y_2=2\sigma+\kappa+\rho+\max(2\sigma,y_0)$ there is a source
depth-one class $a$ and coefficients $o_b\in O$ such that
\[
                p^{y_2} w=F a+\sum_b o_b y_b
                \quad\text{in the target depth-one quotient}.
\]
Indeed multiply the preceding lifting relation by the remaining
nonnegative power of $p$.

By P81F6, $p^\sigma a$ lifts to depth $M$ on the source and
every $p^\sigma y_b$ has a fixed choice of lift at depth $M$ on
the target. The depth-$M$ map to depth one is injective, so
the multiplied equality already holds between these lifted
depth-$M$ classes. Apply the finite Herr class map of P81F6.
Choose one integral finite cocycle for each target basis lift.
The difference between $p^{\sigma+y_2}$ times a tested target
cocycle and this combination belongs to the image of an integral
finite source cocycle modulo an integral target coboundary.
This is precisely membership in $F.\mathrm{ImF}$ plus the span
of the chosen at most $6m$ exceptions. The statement passes to
the entire $J_c$ by linearity. This proves the second scaled
estimate of A81R4 with no dependence on the number of tested
cocycles.
% END U81L1

% BEGIN U81L2
\paragraph{U81L2. The normalized row length at constant 160.}
Let $F:T_1\twoheadrightarrow T_2$ be any map in the original
$\mathrm{RowAt}$ at a twisted ramified index and a good point.
H80M1--H80M2 give the crystalline and numerical hypotheses used
in A81R1--A81R5 on this very map and each good coefficient
extension. A81R11 gives the first scaled inclusion in A81R4 for
every target finite cocycle. U81L1 gives its second inclusion on
each of the $n=[K_0:\mathbb Q_p]$ letters. B81L5 and A81R2 give
the coarse exponent $E$ for all target classes, including the
exceptions. Thus the length proof in A81R5 applies without an
extra local lifting premise. Its two quotients are killed by
$p^Y$ and $p^E$ and have at most $(2n+2)r$ and $6nm$ generators
respectively. Dividing their combined length by the length of
$O/(p)$ gives the asserted estimate with constant $160$.
In particular every map is covered by the rational ceiling
$B_w$ displayed in H80M2. This proves the original
$\mathrm{LengthBoundAtPlaces}$ for the all-row declaration on
the chosen hub. It includes every graph lattice of I81G3 and the
actual matrix-lattice target at $w$. The conclusion here uses
the signed windows determined by the two distinguished primes.
% END U81L2
